%% SpecExec: Methodology
%% Speculative Execution Vulnerability Detection via Graph Neural Networks
%% Compile with: pdflatex specexec_methodology.tex

\documentclass[10pt,twocolumn]{article}

\usepackage[margin=1in]{geometry}
\usepackage{amsmath,amssymb}
\usepackage{booktabs}
\usepackage{graphicx}
\usepackage{hyperref}
\usepackage{xcolor}
\usepackage{listings}
\usepackage{tabularx}
\usepackage{multirow}
\usepackage{array}
\usepackage{paralist}
\usepackage{enumitem}
\usepackage{microtype}
\usepackage{caption}
\usepackage{subcaption}

\hypersetup{
  colorlinks=true,
  linkcolor=blue!60!black,
  citecolor=green!40!black,
  urlcolor=blue!60!black,
}

\lstset{
  basicstyle=\ttfamily\small,
  frame=single,
  xleftmargin=0.5em,
  breaklines=true,
  columns=fullflexible,
}

\title{\textbf{SpecExec: Automated Detection of Speculative Execution Vulnerabilities}}

% \author{
%   Ritvik Gupta \\
%   University of Edinburgh \\
%   \texttt{ritvikg2@andrew.cmu.edu}
% }

\date{}

\begin{document}

\maketitle

\begin{abstract}
We present \textsc{SpecExec}, a system for automated \emph{triage} of
speculative-execution gadget candidates in compiled assembly code. Because
exploitability is a microarchitecture-relative property, we treat the static
label as a candidate-ranking signal for downstream verification rather than a
standalone vulnerability verdict. The system classifies instruction windows into
nine classes---\textsc{Spectre~V1},
\textsc{V2}, \textsc{V4}, \textsc{L1TF}, \textsc{MDS}, \textsc{Retbleed},
\textsc{Inception}, \textsc{Branch History Injection} (BHI), and
\textsc{Benign}---without access to source code, compiler metadata, or
architectural ground truth. Our approach combines Program Dependence Graph
(PDG) representation with learned edge-type attention, 58-dimensional
handcrafted inline features encoding attack-class-specific primitives, and
provably correct data augmentation. On a locked test set derived from a
multi-source corpus (3~architectures, 4~compilers, 5~optimization levels),
our best model achieves \textbf{95.45\%} accuracy with a val/test gap of
under 1\,pp, substantially outperforming prior pattern-matching and
random-forest baselines. We then move the system's hand-authored,
ISA-specific feature engineering behind a declarative ISA specification that
drives graph construction (verified identical to the hand-built pipeline on
437{,}000 edges, and cross-checked against an \emph{independent}
assembler/disassembler oracle that exposes 274 latent classification bugs the
self-comparison cannot). A self-supervised assembly encoder supplies node
features that, under a multi-seed equivalence test, do \emph{not} replace the
hand-designed node vector (learned is significantly lower and not statistically
equivalent) but are strongly \emph{complementary} at the sequence level
($+10.6$\,pp macro-F1). A leakage-controlled evaluation (group- and
optimization-level hold-out) further shows the random-split numbers are
optimistic, with minority-class macro-F1 dropping sharply under a held-out
optimization level---an honest bound on generalization. Finally, the same
representation supports a class- and architecture-conditioned generator that
synthesizes novel, grammar-valid gadgets ($5.7\times$ class-conditioning lift,
$\geq 96\%$ ISA-purity)---the front half of a pipeline aimed at automated
discovery of minimal speculative-leakage gadgets across instruction sets.
\end{abstract}

\section{Problem Formulation}
\label{sec:problem}

Let $S = (i_1, i_2, \ldots, i_n)$ be a \emph{window}---a contiguous
sequence of assembly instructions extracted from a compiled binary. We seek a
classifier $f: S \mapsto \mathcal{C}$ where
\begin{align*}
  \mathcal{C} = \{&\text{BENIGN}, \text{SPECTRE\_V1}, \text{V2}, \text{V4},\\
                  &\text{L1TF}, \text{MDS}, \text{RETBLEED},
                   \text{INCEPTION}, \text{BHI}\}.
\end{align*}

The task is strictly \emph{assembly-level}: no source code, no compiler IR,
no runtime information. Vulnerability presence is determined by structural
patterns in the instruction sequence---existence of a speculation gadget, a
secret-dependent access, and a covert channel transmitter---all of which must
be detectable from register/memory dependency structure and instruction
semantics without symbolic execution.

The nine classes reflect the primary microarchitectural covert-channel
families (Table~\ref{tab:classes}). BHI (Branch History Injection) was added
in this work to distinguish it from Spectre~V2; the distinction lies in the
exploitation vector (indirect branch prediction vs.\ branch history table
poisoning) and manifests in different call patterns and gadget structure at
the assembly level.

\begin{table}[t]
\centering
\caption{Vulnerability classes and distinguishing assembly primitives.}
\label{tab:classes}
\small
\begin{tabularx}{\columnwidth}{@{}lX@{}}
\toprule
Class & Key assembly primitive(s) \\
\midrule
SPECTRE\_V1   & Bounds-bypass: \texttt{cmp; j*; ldr [base+idx*s]} \\
SPECTRE\_V2   & Indirect branch: \texttt{blr xN / jmp *\%rN} \\
SPECTRE\_V4   & Store-bypass: \texttt{str; (spec) ldr} same addr \\
L1TF          & Page-probe: \texttt{shl \$12; movq (\%r1,\%r2)} \\
MDS           & \texttt{movntdqa; verw; mfence} (MFBDS/MFPDS) \\
RETBLEED      & Return gadget: \texttt{push; nop*; pop; ret} \\
INCEPTION     & RSB stuffing: dense \texttt{call/ret} pairs \\
BHI           & \texttt{blr/jmp*} + \texttt{rdtsc+clflush} timing \\
BENIGN        & No speculative transmitter \\
\bottomrule
\end{tabularx}
\end{table}

\section{Data Sources and Collection}
\label{sec:data}

The training corpus draws from four complementary sources to maximize
coverage of compiler fingerprints, instruction set architectures, and
exploit styles.

\subsection{Internal Handcrafted Vulnerabilities}
The primary source of labeled attack-class data is a set of hand-authored C
files encoding canonical vulnerability patterns for each class (directory
\texttt{c\_vulns/}). Each C file is compiled with GCC and Clang across
optimization levels O0--O3 for x86\_64 and ARM64, yielding diverse assembly
realizations of the same logical gadget. The compiler invocations
intentionally include platform-specific flags (\texttt{-march=native},
\texttt{-mretpoline}, etc.) to expose the model to mitigation-related
instruction sequences.

\subsection{Phase~9 Compiler-Synthesized Samples}
A second corpus (tagged \texttt{p9\_*} in the dataset) was produced by a
separate pipeline that compiled known-vulnerable kernel and library code
(Linux kernel PoCs, Spectre PoC repositories) using a broader cross-compiler
matrix: GCC and Clang for x86\_64-linux-gnu and aarch64-linux-gnu at O0 and
O1. These samples provide additional diversity in calling convention,
function prologue/epilogue structure, and inline assembly style not captured
by the handcrafted files.

\subsection{External Verified Sources}
Three external corpora provide independently verified attack samples:

\paragraph{Google SafeSide~\cite{safeside}}
A collection of hardware-confirmed proof-of-concept exploits covering
Spectre~V2 (BTB poisoning), Spectre~V4 (speculative store bypass),
Spectre~RSB (return stack buffer), and partial L1TF scenarios. SafeSide is
distributed under BSD-3-Clause and GPLv2; we compiled each source file with
Apple Clang for both \texttt{x86\_64-apple-macos} and
\texttt{arm64-apple-macos} at O0--O3 (200 total samples).

\paragraph{Spectector Benchmarks~\cite{spectector}}
Formally verified Spectre-V1 assembly files from three real compilers (Intel
ICC, GCC, Clang) applied to 15 Kocher variant programs~\cite{kocher1998}.
We include only the ``vanilla'' (unpatched, non-fenced) variants to prevent
label contamination from lfence-inserted mitigations. These provide diversity
in V1 gadget structure beyond our internal clang-only corpus (101 samples).

\paragraph{FastSpec~\cite{fastspec}}
A dataset of 956K synthetic Spectre-V1 gadgets generated by formal
synthesis. We sample up to 300 records to maintain class-ratio balance;
FastSpec records are used for training only and excluded from the validation
split (``train-only'' partition) because their synthetic distribution may not
represent the natural test distribution.

\subsection{Benign Code}
Benign (SPECTRE-free) assembly is collected from 1,031 GitHub repositories
spanning diverse application domains (networking, compression, text
processing) compiled with GCC and Clang for ARM64 at O1--O3. The breadth of
benign sources is essential: \textsc{SpecExec} must generalize to arbitrary
production code rather than memorizing the negative class.

\subsection{Dataset Statistics}

Table~\ref{tab:data} summarizes the final training and test splits. The test
set is \emph{locked} from v53 of the pipeline and has not been modified
across subsequent versions, enabling direct comparison of model improvements.

\begin{table*}[t]
\centering
\caption{v55 dataset statistics (training and locked test set).}
\label{tab:data}
\small
\begin{tabular}{@{}lrrr@{}}
\toprule
Class & Train & Test & Total \\
\midrule
BENIGN                      & 2840 & 1031 & 3871 \\
BRANCH\_HISTORY\_INJECTION  &  318 &   67 &  385 \\
INCEPTION                   &  518 &   94 &  612 \\
L1TF                        &  582 &   37 &  619 \\
MDS                         &  445 &   45 &  490 \\
RETBLEED                    &  298 &   75 &  373 \\
SPECTRE\_RSB                &  470 &    1 &  471 \\
SPECTRE\_V1                 &  644 &   42 &  686 \\
SPECTRE\_V2                 &  485 &  154 &  639 \\
SPECTRE\_V4                 &  347 &  124 &  471 \\
\midrule
\textbf{Total}              & \textbf{6947} & \textbf{1670} & \textbf{8617} \\
\bottomrule
\end{tabular}
\end{table*}

\section{Data Engineering}
\label{sec:engineering}

\subsection{Assembly Preprocessing}
Raw assembly files are parsed with a custom extractor that discards
assembler directives, debug metadata, and section headers. Instruction
windows of 8--64 instructions are extracted around detected vulnerability
triggers (conditional branches, indirect branches, return instructions).

\paragraph{Function-name neutralization.}
Call targets, \texttt{adrp}/\texttt{:lo12:} address-load pairs (ARM64), and
\texttt{leaq}~\texttt{foo(\%rip)} sequences (x86-64) expose function names
as plaintext tokens. Without neutralization, a model trained on
\texttt{callq~flush\_reload} would learn to classify by the callee name
rather than structural properties. All such targets are replaced with the
literal token \texttt{<fn>} before any feature extraction.

\paragraph{Deduplication.}
Each instruction sequence is SHA-256 hashed after preprocessing. The test
set's hashes are added to a \emph{forbidden set} before any training-set
construction; any training record whose hash matches a test record is
silently dropped. This prevents train-test contamination across pipeline
versions.

\subsection{Label Quality Control}
During construction of the v54 corpus, a BHI label audit revealed that 517
records tagged as BHI originated from OpenSSL library code (functions
\texttt{p15}, \texttt{p19}) that lacks any speculation transmitter. These
were removed, leaving only structurally confirmed BHI gadgets. This audit was
guided by the inline feature vector: sequences flagged as BHI but scoring
zero on all indirect-branch and cache-probe inline features were prioritized
for review.

\subsection{Class-Proportional Augmentation}
\label{sec:augmentation}

The raw corpus is class-imbalanced: the benign class dominates by a factor
of~$\sim$10$\times$ over rare attack classes such as L1TF and RETBLEED. To
address this without na\"ive oversampling (which would duplicate identical
sequences), we apply nine structure-preserving augmentation transforms.

\paragraph{Provability requirement.}
Each augmentation must preserve the \emph{attack validity invariant}: the
augmented sequence must exhibit the same vulnerability class as the original.
We formalize this as: (1)~the speculation trigger must remain present, (2)~no
new covert channel or secret access must be introduced, and (3)~mitigating
instructions (fences, \texttt{verw}) must not be weakened or removed.

We verified all nine transforms via a 143-test suite
(\texttt{tests/augmentation/}). The test suite identified six bugs in the
prior implementation (Table~\ref{tab:aug-bugs}); all were fixed before
training the v55 model.

\begin{table}[t]
\centering
\caption{Augmentation bugs found and fixed by the provability test suite.}
\label{tab:aug-bugs}
\footnotesize
\begin{tabularx}{\columnwidth}{@{}p{1.8cm}X@{}}
\toprule
Function & Bug \& Fix \\
\midrule
\texttt{rename\_regs} & Pool collision: \texttt{x86\_r64} pool contained
  \texttt{r8--r15} (belonging to \texttt{x86\_r}), so \texttt{rcx} and
  \texttt{r9} could both map to \texttt{r13}. Fixed: disjoint pools;
  \texttt{X86\_REG} extended to match \texttt{rax/rbx/rcx/rdx/rsi/rdi}. \\
\addlinespace
\texttt{can\_swap} & x86 branches (\texttt{jne}, \texttt{jge}) not
  guarded; reordering across conditional jumps allowed. Fixed: added
  \texttt{X86\_BRANCH\_COND} check symmetric to existing ARM64 guard. \\
\addlinespace
\texttt{FLAG\_CLOBBER} & Word-boundary regex missed size-suffixed opcodes
  (\texttt{addq}, \texttt{subl}). Fixed: optional suffix \texttt{[qlwb]?}. \\
\addlinespace
\texttt{BARRIER\_SYN.} & \texttt{dsb sy} entry listed weaker variants
  (\texttt{dsb ish}, \texttt{dsb ishst}), violating ``only upgrade'' contract.
  Fixed: each entry lists only equal-or-stronger barriers. \\
\addlinespace
\texttt{X86\_LOAD} & Bare \texttt{mov}/\texttt{lea} missed
  \texttt{movq}/\texttt{movl}/etc.\ in barrier-counterfactual insertion
  on AT\&T size-suffixed x86. Fixed: optional size suffix. \\
\addlinespace
\texttt{flip\_branch} & Guarded \texttt{jmp~*}/\texttt{call~*} but missed
  AT\&T 64-bit forms \texttt{jmpq~*}/\texttt{callq~*}. Fixed: guard
  extended to cover both forms. \\
\bottomrule
\end{tabularx}
\end{table}

\paragraph{The nine transforms and their validity invariants.}

\begin{enumerate}[leftmargin=1.5em,itemsep=1pt]

\item \textbf{Register rename.}
A bijective, family-scoped renaming drawn from a shuffled pool of
callee-saved registers. ABI-special registers (\texttt{sp}, \texttt{xzr},
\texttt{rsp}, \texttt{rbp}) are preserved. Validity: instruction semantics
are register-independent; the gadget structure (def-use chains,
branch→load ordering) is unchanged.

\item \textbf{NOP insertion.}
Inserts \texttt{nop} instructions with probability 0.15, except within a
$\pm 2$-instruction window around conditional branches or memory loads
(speculation-critical positions). Validity: NOPs do not affect functional
semantics; the critical speculation window is not disrupted.

\item \textbf{Local instruction swap.}
Swaps adjacent instruction pairs when \texttt{can\_swap}$(a, b) = \text{True}$:
no register def-use dependence, no fence, no branch in either slot.
Validity: independence is tested by def-use analysis; swapping independent
instructions does not change the program's observable behavior or the
gadget's speculation trigger.

\item \textbf{Immediate perturbation.}
Adds $\pm$1--3 to non-critical immediates. A set $\mathcal{C}$ of
attack-semantically meaningful constants is never modified:
$\mathcal{C} = \{0, 1, -1, 4096, 64, 128, 256, 255, 8, 16, 32, \ldots\}$.
This set includes cache-line size (64), page size (4096), and
common secret-byte iteration bounds (255). Validity: non-critical immediates
are loop indices, offsets, or register sizes that do not affect the
classification of the gadget.

\item \textbf{Stride synonym swap.}
Replaces numeric literals with hexadecimal equivalents (e.g.,
\texttt{\$4096~$\leftrightarrow$~\$0x1000}). Value-preserving by
construction; the test suite verifies numeric equality after swap.

\item \textbf{Equivalent idiom substitution.}
Replaces flag-writing zero idioms (\texttt{xor \%rax, \%rax}
$\to$ \texttt{movq \$0, \%rax}) only when no downstream flag consumer
(\texttt{jcc}/\texttt{setcc}/\texttt{cmovcc}) appears before the next
flag-clobbering instruction. ARM substitutions are unconditionally
flag-neutral (\texttt{mov x0, \#0}~$\leftrightarrow$~\texttt{eor x0, x0, x0}).

\item \textbf{Branch polarity flip.}
Inverts the first conditional branch (\texttt{jge~$\to$~jl}, \texttt{b.lt~$\to$~b.ge}).
Refused for sequences containing \texttt{ret}, \texttt{br}, \texttt{blr},
\texttt{jmpq~*}, or \texttt{callq~*} (return-based gadgets). Validity:
the Spectre-V1 gadget triggers on the speculated path regardless of branch
polarity; flipping the taken/not-taken direction does not remove the gadget.

\item \textbf{Barrier variant swap.}
Replaces a memory barrier with an equal-or-stronger variant (e.g.,
\texttt{lfence~$\to$~mfence}, \texttt{dsb~ish~$\to$~dsb~sy}). Downgrades
are explicitly forbidden. Validity: a stronger barrier stops at least as
much speculative execution as the original, so the security label is preserved.

\item \textbf{Slice recomposition.}
Splits the window into three contiguous chunks and permutes their order when
(a)~no chunk contains control flow, and (b)~the reordered sequence's
live-in register set does not grow relative to the original. Validity:
the live-set walk catches free-use violations that would arise from reordering
past a definition.

\end{enumerate}

For each original record, the nine transforms are applied independently and
capped per class (BENIGN: $\leq$2500 augmented; attack classes: $\leq$600--800).
Records from the phase~9 compiler pipeline (tagged \texttt{p9\_*}) and the
external corpora are treated as originals and \emph{not} re-augmented, since
their diversity is already structural (different compilers, optimization
levels).

\section{Feature Engineering}
\label{sec:features}

\subsection{Program Dependence Graph Construction}
\label{sec:pdg}

Each instruction window is converted to a Program Dependence Graph (PDG) $G
= (V, E)$ where nodes $V$ represent instructions and edges $E$ encode
dependence and speculation-specific structural relationships.

\paragraph{Node features (41 dimensions).}
Each node $v_i$ is represented by a 41-dimensional feature vector:
\begin{itemize}[itemsep=0pt]
  \item \textbf{Opcode category one-hot (19 dims):} LOAD, STORE,
    BRANCH\_COND, BRANCH\_UNCOND, CALL, CALL\_INDIRECT, RET,
    JUMP\_INDIRECT, COMPARE, ARITHMETIC, LOGIC, SHIFT, FENCE,
    CACHE, TIMING, MOVE, STACK, NOP, OTHER.
  \item \textbf{Memory access type one-hot (5 dims):} NONE, STACK, HEAP,
    INDEXED, INDIRECT.
  \item \textbf{Register degree features (2 dims):} number of destination
    and source registers, normalized to [0, 1].
  \item \textbf{Speculation flags (14 dims, binary):}
    \texttt{is\_serializing}, \texttt{is\_cache\_probe}, \texttt{is\_branch},
    \texttt{is\_indirect\_branch}, \texttt{is\_memory\_access},
    \texttt{is\_timing\_source}, \texttt{is\_secret\_source},
    \texttt{is\_transmitter}, \texttt{is\_lfence}, \texttt{is\_mfence\_or\_sfence},
    \texttt{is\_verw}, \texttt{is\_prefetch}, \texttt{is\_nontemp\_load},
    \texttt{is\_gather}.
\end{itemize}

\paragraph{Edge types (9 types).}
\begin{enumerate}[itemsep=0pt]
  \item \textbf{DATA\_DEP} --- register def-use dependency.
  \item \textbf{CONTROL\_FLOW} --- sequential instruction ordering.
  \item \textbf{SPEC\_CONDITIONAL} --- the speculation edge from a conditional
    branch to each instruction on its predicted-taken path, encoding the
    transient execution window.
  \item \textbf{SPEC\_INDIRECT} --- speculation from an indirect branch/call.
  \item \textbf{SPEC\_RETURN} --- speculation from a return instruction
    (RSB-based).
  \item \textbf{MEMORY\_ORDER} --- store-load ordering constraints within the
    window.
  \item \textbf{CACHE\_TEMPORAL} --- \texttt{clflush}/\texttt{prfm} to
    subsequent load of the same cache set.
  \item \textbf{FENCE\_BOUNDARY} --- fence instruction to all instructions
    that depend on speculation it serializes.
  \item \textbf{RSB\_CHAIN} --- call to return pairs within a 15-instruction
    window, encoding RSB stuffing gadgets (dense in INCEPTION; sparse in
    RETBLEED).
\end{enumerate}

Edge types 3--5 capture the semantic structure of speculative execution
directly in the graph topology. The SPEC\_CONDITIONAL edge, for instance,
creates a path from the branch to the secret-dependent load that would not
appear in a standard data-flow graph; this path is the structural signature
the model learns to detect. The RSB\_CHAIN edge was added to separate
INCEPTION (RSB stuffing: many call/ret pairs) from RETBLEED (single return
gadget): prior versions of the model confused these classes due to their
shared use of \texttt{ret} instructions.

\subsection{Inline Handcrafted Features (58 dimensions)}
\label{sec:inline}

In addition to the graph representation, we compute a 58-dimensional
handcrafted feature vector from the raw instruction sequence. These features
encode domain knowledge about each vulnerability class that the GNN may take
many training samples to discover independently.

\paragraph{Group~1: Opcode frequency fractions (28 dims).}
Fraction of instructions matching each of 28 security-relevant opcodes:
\texttt{nop}, \texttt{ret}, \texttt{retq}, \texttt{call}, \texttt{callq},
\texttt{bl}, \texttt{blr}, \texttt{br}, \texttt{cmp}, \texttt{cmn},
\texttt{test}, \texttt{tst}, \texttt{ldr}, \texttt{ldp}, \texttt{str},
\texttt{stp}, \texttt{movq}, \texttt{movl}, \texttt{clflush},
\texttt{clflushopt}, \texttt{rdtsc}, \texttt{rdtscp}, \texttt{lfence},
\texttt{mfence}, \texttt{sfence}, \texttt{verw}, \texttt{movntdqa},
\texttt{adrp}. Notably, \texttt{endbr64} was excluded after ablation showed
it amplified RETBLEED confusion and degraded L1TF recall.

\paragraph{Group~2: Structural ratios (6 dims).}
Fraction of x86-only instructions, fraction of ARM-only instructions,
fraction of indirect operations, fraction of branches, loads, and stores.
These distinguish architecture families and control vs.\ data-heavy gadgets.

\paragraph{Group~3: Structural pattern counts, normalized (7 dims).}
Maximum NOP run length (normalized by sequence length), call/return pair
density, indexed-load density, \texttt{clflush}+load pattern density,
branch-after-compare density, indirect-then-load density, unique-opcode
fraction. These capture the high-level ``shape'' of a gadget.

\paragraph{Group~4: Speculative primitive presence flags (10 dims, binary).}
\texttt{has\_rdtsc}, \texttt{has\_verw}, \texttt{has\_movntdqa},
\texttt{has\_lfence}, \texttt{has\_clflush}, \texttt{has\_indirect},
\texttt{has\_nop\_run\_3plus}, \texttt{has\_indexed\_load},
\texttt{has\_call\_ret\_pair}, \texttt{has\_clflush\_load}. These are
binary indicators of the presence of attack-primitive instructions.

\paragraph{Group~5: Ratio features (4 dims).}
\texttt{ret\_call\_ratio}, \texttt{nop\_ret\_ratio},
\texttt{indirect\_cond\_ratio}, \texttt{load\_store\_ratio}. The
\texttt{ret\_call\_ratio} is particularly discriminating: INCEPTION uses
many call/ret pairs (ratio $\approx$ 1) while RETBLEED has an unmatched
\texttt{ret} (ratio $> 1$).

\paragraph{Group~6: Caller-context (1 dim).}
Binary flag indicating whether any call target in the sequence matches a
regex of attack-function names (\texttt{flush\_reload}, \texttt{gadget\_*},
\texttt{victim\_function}, etc.). This feature fires on research-PoC code
that explicitly names its components and is always zero on production code,
providing a useful signal without causing label leakage.

\paragraph{Group~7: Meltdown/L1TF page-probe (2 dims).}
Two features specific to the L1TF and Meltdown covert channel. The
\emph{page-probe transmitter} in Meltdown-based attacks uses a page-size
shift to index into a probe array: \texttt{shlq \$12, \%rax} followed within
four instructions by an indexed load \texttt{movq (\%base,\%rax,1), \%rdi}.
The shift amount $2^{12} = 4096$ bytes equals the x86 page size; this ensures
each byte value of the secret indexes a distinct page, preventing cache
set aliasing. Feature \texttt{has\_page\_probe\_load} (binary) fires on any
such pattern; \texttt{page\_probe\_density} gives the normalized count.

Key design decisions in this feature:
(a)~Only shifts \$9--\$15 are page-probes; \$6 is a cache-line probe
(stride $64B = 2^6$) shared by MDS and L1TF, hence non-discriminating.
(b)~AT\&T indexed stores (\texttt{movb \%al, (\%rsi,\%rcx)}) have the indexed
operand as the \emph{destination} and no trailing \texttt{,~\%reg}---the regex
requires a trailing destination register to distinguish loads from stores.
(c)~Clang omits the scale-1 suffix in two-register addressing modes
(\texttt{(\%r14,\%rax)} vs.\ GCC's \texttt{(\%rsi,\%rax,1)}); the regex
matches both forms.

\subsection{Global Graph Features (5 dims)}

Five instruction-fraction statistics computed from the raw sequence and
provided as a separate branch:
\begin{enumerate}[label=(\arabic*)]
  \item NOP fraction,
  \item indirect instruction fraction,
  \item return instruction fraction,
  \item \texttt{verw} fraction,
  \item \texttt{movntdqa} fraction.
\end{enumerate}
These are computed identically for train and test from their own sequences
with no training-set normalization, eliminating any leakage.

\subsection{Architecture Embedding (8 dims)}

A learned 8-dimensional embedding indexed by the assembly architecture
(\texttt{x86\_64}, \texttt{arm64}, \texttt{riscv}, etc.). This accounts for
the systematic differences in opcode vocabulary and addressing syntax across
ISAs without hard-coding ISA-specific rules.

\section{Model Architecture}
\label{sec:model}

\subsection{Graph Isomorphism Network with Edge Features (GINE)}

We use the Graph Isomorphism Network with Edge features
(GINE)~\cite{hu2020strategies} as our backbone. GINE generalizes GIN's
sum-aggregation to incorporate edge features:
\begin{equation*}
\begin{split}
  h_v^{(k)} &= \text{MLP}^{(k)}\!\biggl(
    (1 + \varepsilon^{(k)}) \cdot h_v^{(k-1)} \\
    &\quad + \!\!\sum_{u \in \mathcal{N}(v)}\!\! \text{ReLU}\!\bigl(
      h_u^{(k-1)} + e_{uv}
    \bigr)\biggr),
\end{split}
\end{equation*}
where $e_{uv}$ is the edge-type embedding for edge $(u, v)$ and
$\varepsilon^{(k)}$ is a learned scalar. GNN architectures that ignore edge
types (standard GCN, GAT) cannot distinguish, e.g., a DATA\_DEP edge from a
SPEC\_CONDITIONAL edge---a distinction that is the \emph{defining} structural
property of the vulnerability.

\paragraph{Layer configuration.}
We use $K = 3$ GINE layers with hidden dimension $d = 128$, dropout 0.5,
and batch normalization after each layer. Residual connections
$h^{(k)} \leftarrow \text{LayerNorm}(h^{(k-1)} + h_\text{new}^{(k)})$
prevent vanishing gradients across layers.

\paragraph{Edge-type learnable scaling.}
Each edge type $t$ is associated with a learnable scale $s_t \in \mathbb{R}_+$:
the edge embedding for type $t$ is multiplied by $s_t$ before message
passing. This allows the model to weight different structural relationships
differently. In converged models, SPEC\_CONDITIONAL and SPEC\_RETURN
consistently scale up ($s \approx 1.1$--1.2); CONTROL\_FLOW consistently
scales down ($s \approx 0.75$--0.88), reflecting that raw sequential order
is less informative than speculative dependency structure.

\paragraph{Virtual node.}
A learnable global state vector $\mathbf{v} \in \mathbb{R}^d$ is maintained
across layers and broadcast to all nodes before each GINE update. The
broadcast is gated by a learned scalar initialized to $\sigma(-2) \approx
0.12$ to prevent the virtual node from dominating local message passing in
early training. The virtual node enables global information exchange between
nodes that would otherwise be many hops apart, addressing the limited
receptive field of 3-layer GNNs.

\paragraph{JK-Cat readout.}
Jumping Knowledge Concatenation~\cite{xu2018jk} aggregates the hidden states
from all layers plus the initial encoding: $\mathbf{h}_v^{\text{JK}} =
[\mathbf{h}_v^{(0)} \| \mathbf{h}_v^{(1)} \| \mathbf{h}_v^{(2)} \|
\mathbf{h}_v^{(3)}]$, giving a 512-dimensional representation per node before
readout. This preserves both local (single-hop) and global (3-hop) structure.

\paragraph{Attention readout.}
An attention network $\alpha_v = \text{softmax}(\mathbf{w}^\top
\tanh(\mathbf{W} \mathbf{h}_v^{\text{JK}}))$ weights each node's contribution
to the graph-level embedding. Security-relevant nodes (memory accesses,
branches, fences) receive higher attention weights than boilerplate
instructions (stack frame setup, function calls unrelated to the gadget).
The pooled graph representation $\mathbf{g} = \sum_v \alpha_v
\mathbf{h}_v^{\text{JK}}$ is projected to a 256-dimensional fusion vector.

\subsection{Feature Fusion and Classification}

The inline handcrafted features (58-dim) are encoded by a two-layer MLP
(Linear-BN-ReLU-Dropout-Linear-BN-ReLU-Dropout) to a 256-dim vector. The
global graph features (5-dim) are projected to 32 dims; the architecture
embedding provides 8 dims. The four representations are concatenated to form
a $256 + 256 + 32 + 8 = 552$-dimensional combined vector, then classified by
a two-layer MLP with 128 hidden units:
\[
  \hat{y} = \text{MLP}_\text{cls}\!\left(
    [\mathbf{g}_\text{graph} \| \mathbf{g}_\text{inline} \|
     \mathbf{g}_\text{global} \| \mathbf{g}_\text{arch}]
  \right).
\]

An auxiliary head $\hat{y}_\text{aux} = \mathbf{W}_\text{aux}
\mathbf{g}_\text{inline}$ predicts the label from inline features alone.
The auxiliary loss ($\lambda_\text{aux} = 0.3$) regularizes the inline
feature branch and ensures the graph branch cannot simply delegate all
classification to it.

\paragraph{Total parameters: 618,059.}

\section{Training Protocol}
\label{sec:training}

\subsection{Validation Split}

The training set is divided into a model-selection validation split using
\texttt{StratifiedGroupKFold} with 10 splits. The ``group'' field encodes
the source file of each record; no group appears in both the training and
validation partitions. This prevents the model from exploiting repeated
windows from the same file to inflate validation accuracy while generalizing
poorly to new files. FastSpec records (synthetically generated, different
distributional properties) are excluded from the validation pool.

\subsection{Loss Function}

The total training loss is:
\[
  \mathcal{L} = \mathcal{L}_\text{CE} + \lambda_\text{con}
    \mathcal{L}_\text{SupCon} + 0.3\, \mathcal{L}_\text{aux},
\]
where $\mathcal{L}_\text{CE}$ is class-weighted cross-entropy,
$\mathcal{L}_\text{SupCon}$ is supervised contrastive
loss~\cite{khosla2020supervised}, and $\mathcal{L}_\text{aux}$ is the
auxiliary inline-features cross-entropy head.

\paragraph{Class-weighted CE.}
Class weights are $w_c = 1/\sqrt{n_c}$ normalized to unit mean. The
square-root scaling is less aggressive than $1/n_c$ (which can destabilize
training for very rare classes like SPECTRE\_RSB with $n = 3$ originals)
while still providing meaningful rebalancing.

\paragraph{Supervised contrastive loss.}
The projection head maps the combined 552-dim representation to a 128-dim
unit sphere. The SupCon loss pulls together examples of the same class while
pushing apart examples of different classes~\cite{khosla2020supervised}.
Thirteen \emph{hard-negative pairs}---class pairs that are structurally
similar but semantically distinct---receive an additional multiplicative
penalty of 2.0 (Table~\ref{tab:hard-neg}). These pairs were identified from
confusion matrix analysis across prior model versions.

\begin{table}[t]
\centering
\caption{Hard-negative class pairs receiving 2$\times$ contrastive penalty.}
\label{tab:hard-neg}
\footnotesize
\begin{tabularx}{\columnwidth}{@{}lX@{}}
\toprule
Pair & Structural similarity \\
\midrule
(V1, BHI)             & Cond.\ branch + indexed load \\
(INCEPTION, BHI)      & RSB-adjacent call/ret patterns \\
(MDS, RETBLEED)       & x86 timing helpers structurally identical \\
(INCEPTION, RETBLEED) & RSB exploitation \\
(L1TF, V1)            & L1 cache side-channel \\
(L1TF, V4)            & Memory-access-based gadget \\
(MDS, V4)             & Memory ordering dependence \\
(V1, V4)              & Branch bypass vs.\ store bypass \\
(V2, BHI)             & Indirect branch manipulation \\
(V2, INCEPTION)       & Indirect branch vs.\ RSB \\
(RETBLEED, INCEPTION) & Return-based exploitation \\
(MDS, L1TF)           & \texttt{clflush}+load timing loop \\
(V2, RETBLEED)        & Indirect vs.\ return speculation \\
\bottomrule
\end{tabularx}
\end{table}

\paragraph{Contrastive warmup.}
$\lambda_\text{con}$ is linearly warmed up from 0 to its target value (0.5)
over the first 10 epochs. This prevents the contrastive loss from dominating
early training before the classifier has established meaningful embeddings.

\subsection{Optimizer and Scheduler}

AdamW~\cite{loshchilov2019decoupled} with initial learning rate $10^{-3}$
and weight decay $5 \times 10^{-4}$, batch size 32. A cosine annealing
schedule reduces the learning rate to near zero over 120 epochs. Early
stopping on validation accuracy with patience 15: training terminates when
validation accuracy does not improve for 15 consecutive epochs.

\subsection{Hardware}

All models are trained on CPU (Apple M-series, MPS disabled on Python~3.14
due to known stability issues). Training time is approximately 10--13 seconds
per epoch; typical convergence occurs at epochs 12--20.

\section{Results and Interpretation}
\label{sec:results}

\subsection{Overall Performance}

Table~\ref{tab:results-full} gives precision, recall, and F1 for every
class on the locked 1670-record test set; both v54 (pre-bug-fix augmentation)
and v55 (provably correct augmentation) use the identical architecture and
hyperparameters so any difference is attributable solely to data quality.
Figure~\ref{fig:cm} shows the normalized confusion matrices.

\begin{table*}[t]
\centering
\caption{Full per-class metrics on the locked test set ($n_\text{test}=1670$).
  v54: test accuracy 95.75\%, best epoch 19 of 34, val acc 95.92\%.
  v55: test accuracy 95.45\%, best epoch 12 of 27, val acc 94.95\%.
  SPECTRE\_RSB has $n=1$ test sample; its metrics are unreliable.
  Bold marks improvement in recall $\geq$5\,pp; italic marks regression $\geq$5\,pp.}
\label{tab:results-full}
\footnotesize
\begin{tabular}{@{}l rrr rrr r@{}}
\toprule
 & \multicolumn{3}{c}{\textbf{v54}} & \multicolumn{3}{c}{\textbf{v55 (best/3)}} & \\
\cmidrule(lr){2-4}\cmidrule(lr){5-7}
Class & Prec & Rec & F1 & Prec & Rec & F1 & $n$ \\
\midrule
BENIGN
  & 1.000 & 1.000 & 1.000
  & 0.994 & 0.984 & 0.989 & 1031 \\
BRANCH\_HISTORY\_INJ.
  & 0.625 & 0.896 & 0.736
  & 0.756 & 0.881 & 0.814 &   67 \\
INCEPTION
  & 0.945 & 0.915 & 0.930
  & 0.818 & \textit{0.862} & 0.839 &   94 \\
L1TF
  & 0.857 & 0.811 & 0.833
  & 0.958 & \textit{0.622} & 0.754 &   37 \\
MDS
  & 0.894 & 0.933 & 0.913
  & 0.860 & 0.956 & 0.905 &   45 \\
RETBLEED
  & 0.944 & 0.907 & 0.925
  & 0.961 & \textbf{0.973} & 0.967 &   75 \\
SPECTRE\_RSB$^\dagger$
  & 1.000 & 1.000 & 1.000
  & 1.000 & 1.000 & 1.000 &    1 \\
SPECTRE\_V1
  & 0.840 & 1.000 & 0.913
  & 0.808 & 1.000 & 0.894 &   42 \\
SPECTRE\_V2
  & 0.944 & 0.760 & 0.842
  & 0.931 & \textbf{0.877} & 0.903 &  154 \\
SPECTRE\_V4
  & 0.992 & 0.984 & 0.988
  & 0.984 & 0.984 & 0.984 &  124 \\
\midrule
Macro avg
  & 0.904 & 0.920 & 0.908
  & 0.907 & 0.914 & 0.905 & 1670 \\
Weighted avg
  & 0.964 & 0.957 & 0.958
  & 0.958 & 0.954 & 0.955 & 1670 \\
\midrule
\textbf{Overall accuracy}
  & \multicolumn{3}{c}{\textbf{95.75\%}}
  & \multicolumn{3}{c}{\textbf{95.45\%}} & 1670 \\
Val/test gap
  & \multicolumn{3}{c}{0.17\,pp}
  & \multicolumn{3}{c}{$<$1.0\,pp} & \\
\bottomrule
\multicolumn{8}{l}{%
  $^\dagger$ $n=1$ test sample; recall/precision are not statistically meaningful.}
\end{tabular}
\end{table*}

The 0.30\,pp accuracy gap between v54 and v55 is within stochastic training
variance: across three v55 runs, test accuracy ranged 93.41--95.45\%. v54's
anomalously small val/test gap (0.17\,pp) may partly reflect overfitting via
the register-rename pool-collision bug, which inadvertently acted as register
anonymization on some test-adjacent sequences.

\begin{figure*}[t]
\centering
\begin{subfigure}[b]{0.48\textwidth}
  \centering
  \includegraphics[width=\textwidth]{./v54_confusion_matrix.png}
  \caption{v54 --- test accuracy 95.75\%}
  \label{fig:cm-v54}
\end{subfigure}
\hfill
\begin{subfigure}[b]{0.48\textwidth}
  \centering
  \includegraphics[width=\textwidth]{./v55_confusion_matrix.png}
  \caption{v55 --- test accuracy 95.45\% (best of 3 seeds)}
  \label{fig:cm-v55}
\end{subfigure}
\caption{Normalized confusion matrices on the locked test set (1670 records).
  Cell $(i,j)$ shows fraction of true-class-$i$ samples predicted as class $j$.
  Diagonal = recall per class. Off-diagonal dark cells indicate systematic
  confusion: BHI$\to$INCEPTION (indirect vs.\ RSB), L1TF$\to$MDS
  (shared timing loop structure). v55 eliminates the BHI precision deficit
  visible in v54 at the cost of higher L1TF scatter.}
\label{fig:cm}
\end{figure*}

\subsection{Training Convergence}

Figure~\ref{fig:training} shows validation accuracy across epochs for both
versions. v54 requires 34 epochs (best at epoch~19), whereas v55 converges
faster (27 epochs, best at epoch~12), likely because the corrected augmentation
produces a more consistent training signal from epoch~1 onwards.
v55 epoch-1 validation accuracy is 89.7\% vs.\ v54's 75.5\%---the cleaner
augmented data allows faster initial specialization.

\begin{figure*}[t]
\centering
\begin{subfigure}[b]{0.48\textwidth}
  \centering
  \includegraphics[width=\textwidth]{../v54/viz_v54/training_history.png}
  \caption{v54 training history (34 epochs, best at 19)}
  \label{fig:th-v54}
\end{subfigure}
\hfill
\begin{subfigure}[b]{0.48\textwidth}
  \centering
  \includegraphics[width=\textwidth]{../v55/viz_v55/training_history.png}
  \caption{v55 training history (27 epochs, best at 12)}
  \label{fig:th-v55}
\end{subfigure}
\caption{Training convergence. Top: CE loss (solid) and contrastive loss (dashed).
  Bottom: training accuracy (solid) and validation accuracy (dashed). The grey
  vertical line marks the early-stopping checkpoint (best validation accuracy).
  v55 CE loss plateaus higher (0.068 vs.\ 0.040), reflecting a harder training
  set with correctly diverse augmented sequences.}
\label{fig:training}
\end{figure*}

\subsection{Per-Class Interpretation}

\paragraph{Spectre V2 and RETBLEED improvements (+11.7\,pp, +6.6\,pp).}
Both classes benefit from the fixed barrier-variant augmentation.
The prior code allowed \texttt{dsb~sy~$\to$~dsb~ish} downgrades: the model
trained on v54 data saw weakened barriers still tagged as V2/RETBLEED,
introducing confusion with mitigated sequences. With the fix, SPECTRE\_V2
recall rises from 76.0\% to 87.7\% and RETBLEED from 90.7\% to 97.3\%.
BHI precision also improves substantially (0.625~$\to$~0.756) because the
fixed register-rename no longer accidentally conflates indirect-branch patterns
with RSB patterns via register collision.

\paragraph{L1TF regression ($-$18.9\,pp in best v55 run).}
L1TF drops from 81.1\% to 62.2\% recall (though L1TF precision rises:
0.857~$\to$~0.958, indicating fewer false positives). Three contributing
factors: (a)~L1TF has only 84 originals, yielding 464 augmented records
after capping---still sparse relative to MDS (445) despite structural
similarity, making the model sensitive to boundary decisions; (b)~the
now-correct \texttt{swap\_locally} no longer inadvertently crosses the
\texttt{shl/movq} page-probe boundary, eliminating what was an accidental
form of ``hard negative'' augmentation; (c)~high run-to-run variance for
this class ($n_\text{test}=37$): across 3 v55 runs, L1TF recall ranged
62--76\%.

\paragraph{INCEPTION regression ($-$5.3\,pp).}
The fixed register rename removes the x86 pool collision, so INCEPTION samples
no longer silently map \texttt{rcx} and \texttt{r9} to the same output
register. That collapse acted as inadvertent register anonymization; without it
the model must learn genuine structural invariance and needs more training data
to do so reliably.

\subsection{Qualitative PDG Analysis of Confused Pairs}
\label{sec:pdg-analysis}

Figures~\ref{fig:pdg-l1tf-mds}--\ref{fig:pdg-bhi-v2} show the PDGs of
representative test samples from the three most confused class pairs,
rendered with a colorblind-safe palette~\cite{wong2011}.
Node color encodes opcode category; edge color and line style encode the
nine edge types (see legend). Each pair is drawn at the same scale; all
samples are from x86\_64 compiled with Clang.

\begin{figure*}[t]
\centering
\includegraphics[width=\textwidth]{figures/pdg_l1tf_mds.png}
\caption{\textbf{L1TF vs.\ MDS} (structurally similar).
  Both sequences follow the pattern
  \texttt{clflush} (cache/black node) $\to$ \texttt{nop}$^+$ (white) $\to$
  \texttt{shlq \$6} (ALU/gray) $\to$ indexed \texttt{movq} (load/blue),
  connected by DATA\_DEP (sky-blue) and CACHE\_TEMPORAL (black dashed) edges.
  The discriminating structural feature---a \texttt{shlq \$12} page-stride shift
  (L1TF Meltdown transmitter) vs.\ a \texttt{verw} flush node (MDS MFBDS)---is
  absent in these short windows, explaining the 18.9\,pp L1TF recall drop in v55.}
\label{fig:pdg-l1tf-mds}
\end{figure*}

\begin{figure*}[t]
\centering
\includegraphics[width=\textwidth]{figures/pdg_inception_retbleed.png}
\caption{\textbf{INCEPTION vs.\ RETBLEED} (structurally contrasting).
  The INCEPTION PDG (left, 16 nodes, 32 edges) shows a dense subgraph:
  multiple \texttt{callq*} nodes (vermilion, CALL\_INDIRECT) connected by
  RSB\_CHAIN edges (vermilion dash-dot) to three \texttt{retq} nodes
  (vermilion), plus DATA\_DEP chains linking the paired \texttt{leaq}/\texttt{callq}
  instructions.
  The RETBLEED PDG (right, 8 nodes) is a near-empty linear chain: four
  \texttt{nop} nodes (white) with only CONTROL\_FLOW edges (light gray),
  and a single isolated \texttt{retq} node with no RSB\_CHAIN edges.
  RSB\_CHAIN density is the primary discriminating feature;
  RETBLEED recall improves from 90.7\% (v54) to 97.3\% (v55) partly because
  the fixed augmentation no longer produces INCEPTION sequences with
  accidentally collapsed call-register pairs that resemble RETBLEED.}
\label{fig:pdg-inception-retbleed}
\end{figure*}

\begin{figure*}[t]
\centering
\includegraphics[width=\textwidth]{figures/pdg_bhi_v2.png}
\caption{\textbf{BHI vs.\ SPECTRE\_V2} (shared indirect-branch structure).
  Both PDGs are dominated by SPEC\_INDIRECT edges (vermilion dashed) emanating
  from an indirect branch node (\texttt{jmpq*} or \texttt{blr}, vermilion).
  In BHI (left, 8 nodes), a \texttt{movq 8(\%rax)} load precedes the indirect
  jump, forming a two-hop dependency: dispatch-table pointer load $\to$
  SPEC\_INDIRECT $\to$ secret load. In SPECTRE\_V2 (right), SPEC\_INDIRECT
  edges fan directly to all post-branch nodes from a single indirect branch,
  without the preceding table-pointer load chain.
  BHI precision improves from 0.625 (v54) to 0.756 (v55) because the corrected
  register-rename no longer accidentally maps the dispatch-table load register
  to the same destination as the indirect branch target.}
\label{fig:pdg-bhi-v2}
\end{figure*}

\subsection{Stochastic Variance and Reporting}

Run-to-run variance is 1--2\,pp in overall accuracy and up to 15--20\,pp in
per-class recall for small-support classes (L1TF: $n=37$; SPECTRE\_V1: $n=42$;
BHI: $n=67$). For rigorous reporting we recommend: (a)~3--5 seeds, reporting
mean\,$\pm$\,std; (b)~flagging classes with $n_\text{test}<50$ as
statistically unreliable; (c)~treating SPECTRE\_RSB ($n=1$) as unmeasured.

\subsection{Edge Type Scale Analysis}

\begin{table}[t]
\centering
\caption{Converged per-edge-type learnable scales $s_t$ at the best-epoch
  checkpoint. Scale $>1$ means the model amplifies that edge type's signal;
  $<1$ means it suppresses it.}
\label{tab:edge-scales}
\footnotesize
\begin{tabularx}{\columnwidth}{@{}lrrX@{}}
\toprule
Edge type & v54 & v55 & Interpretation \\
\midrule
DATA\_DEP       & 1.093 & 1.003 & Register def-use; moderate signal \\
CONTROL\_FLOW   & 0.743 & 0.747 & Sequential order; suppressed \\
SPEC\_COND.     & 0.982 & 1.149 & Speculative window from branch \\
SPEC\_INDIRECT  & 1.056 & 1.163 & Indirect-branch speculation \\
SPEC\_RETURN    & 1.232 & 1.252 & Return-based speculation; top signal \\
MEMORY\_ORDER   & 1.064 & 1.084 & Store-load ordering \\
CACHE\_TEMPORAL & 1.035 & 1.005 & Flush-then-load temporal ordering \\
FENCE\_BOUNDARY & 0.997 & 1.100 & Serializing fence scope \\
RSB\_CHAIN      & 1.030 & 1.007 & Call-return RSB stuffing pairs \\
\bottomrule
\end{tabularx}
\end{table}

Table~\ref{tab:edge-scales} and Figure~\ref{fig:edge-scales} show the
converged edge-type scales. SPEC\_RETURN is the most amplified type in
both versions ($s=1.23$ and $1.25$ respectively), confirming that
return-based speculative edges are the primary discriminating structural
feature. CONTROL\_FLOW is the most suppressed ($s\approx0.74$ in both),
confirming that raw sequential order is largely redundant given the explicit
dependence and speculative edges. In v55, SPEC\_CONDITIONAL and
SPEC\_INDIRECT scale up noticeably ($+$0.17, $+$0.11 vs.\ v54), consistent
with the improved SPECTRE\_V2 recall: V2 is an indirect-branch attack and
BHI exploits history-table poisoning, both detected via SPEC\_INDIRECT and
SPEC\_CONDITIONAL edges.

\begin{figure*}[t]
\centering
\begin{subfigure}[b]{0.48\textwidth}
  \centering
  \includegraphics[width=\textwidth]{../v54/viz_v54/edge_type_scale_evolution.png}
  \caption{v54 edge scale evolution across training}
\end{subfigure}
\hfill
\begin{subfigure}[b]{0.48\textwidth}
  \centering
  \includegraphics[width=\textwidth]{../v55/viz_v55/edge_type_scale_evolution.png}
  \caption{v55 edge scale evolution across training}
\end{subfigure}
\caption{Per-edge-type learnable scale $s_t$ as a function of training epoch.
  SPEC\_RETURN (solid orange) consistently rises above all others in both
  versions. CONTROL\_FLOW (dashed blue) is the only type that converges
  below 1.0, indicating the model learns to discount raw sequential
  ordering in favor of explicit dependence and speculative structure.}
\label{fig:edge-scales}
\end{figure*}

\section{Automated Feature Engineering and Conditioned Gadget Generation}
\label{sec:specdiscover}

The classifier of Sections~\ref{sec:features}--\ref{sec:results} relies on
hand-authored, ISA-specific feature extractors: opcode taxonomies, AT\&T/ARM64
regular expressions, per-class acceptance filters, and edge-window constants.
Extending the system to new instruction sets or attack classes therefore
requires re-authoring code. This section presents three steps that remove that
manual surface and turn the classifier into the front half of a generative
discovery pipeline whose eventual goal is to synthesize, rank, and
simulator-confirm \emph{minimal} speculative-leakage gadgets across
architectures. We report each step with a verification that holds it to the
existing system as a reference.

\subsection{ISA-Specification--Driven Feature Extraction}
\label{sec:spec-engine}

We replace the hardcoded classification logic in the PDG builder
(Section~\ref{sec:pdg}) with a declarative \emph{ISA specification}: an external
document that lists the opcode-to-category map, an operand/addressing grammar,
speculation-primitive tags, and a pipeline block (speculative / cache / RSB
windows, speculation sources). A generic engine interprets the specification to
classify instructions, so a new ISA is described by a spec file rather than new
Python. Specifications compose by single-level inheritance: a small
\texttt{base} document carries the shared instruction taxonomy, and thin
per-ISA documents (\texttt{x86\_64}, \texttt{arm64}) add pipeline and
realization parameters.

To guarantee the externalized specification does not drift from the reference
system, the \texttt{base} document is exported programmatically from the current
builder (patterns and taxonomies copied verbatim), and two equivalences are
checked on the full corpus (Table~\ref{tab:spec-equiv}): (i) every per-node
decision---opcode category, memory-access type, the 14-dimensional speculation
flags, and register def/use sets---and (ii) the entire constructed graph (node
vectors plus the full multiset of typed edges across all nine edge types). Both
hold with zero mismatches, so the spec-driven builder is a verified drop-in.

Zero mismatches, however, only prove \emph{refactoring fidelity}: the
specification is exported from the builder and checked against that same builder,
so the test is a round-trip and says nothing about whether the classification is
\emph{correct}. To obtain an independent signal we add a second oracle that
shares no code with our regexes: each instruction is assembled with
\texttt{llvm-mc} and its control-flow category is read from the resulting machine
code via \texttt{capstone} instruction groups. Over the $25{,}942$ unique
(arch, instruction) pairs (87.9\% assemble-and-decode coverage), the spec agrees
with this oracle on $98.80\%$ of the covered control-flow instructions, and the
$274$ disagreements are genuine defects the self-comparison cannot surface: (i)
missing mnemonics that silently drop speculation sources to \textsc{other}
(x86 \texttt{jge/jle/jae/jbe}, ARM \texttt{blr} indirect calls, tab-separated
unconditional branches), and (ii) cross-ISA pattern contamination, where x86
register spellings (\texttt{\%bl}, \texttt{\%bpl}) match ARM branch mnemonics in
the merged \texttt{base} grammar. The former are targeted regex fixes; the latter
is the concrete motivation for decomposing \texttt{base} into truly per-ISA
grammars. Both change node categories and therefore require re-export,
re-validation, and retraining, so we report them as findings rather than applying
them silently. The honest correctness statement is thus \emph{98.80\% agreement
with an independent oracle, with 274 triaged defects}---not ``zero mismatches.''

\begin{table}[t]
\centering
\caption{Equivalence of the spec-driven engine to the reference PDG builder,
  on the full v54 corpus. Zero mismatches on every checked quantity.}
\label{tab:spec-equiv}
\footnotesize
\begin{tabularx}{\columnwidth}{@{}Xrr@{}}
\toprule
Check & Units compared & Mismatches \\
\midrule
Node decisions (cat/mem/flags/regs) & 25{,}942 instr. pairs & 0 \\
Full graph, window=20 & 207{,}113 nodes / 437{,}231 edges & 0 \\
Full graph, window=10 & 207{,}113 nodes / 421{,}525 edges & 0 \\
\bottomrule
\end{tabularx}
\end{table}

\subsection{Self-Supervised Learned Node Features}
\label{sec:learned-features}

The 40-dimensional PDG node vector and the 58 inline features encode expert
knowledge that is ISA- and attack-specific. We ask whether a
\emph{self-supervised} encoder can supply the per-node signal with no
hand-design, holding the rest of the pipeline fixed.
Instructions are first normalized to a low-vocabulary form
(registers\,$\to$\,\texttt{<reg>}, immediates\,$\to$\,\texttt{<imm>}, memory
operands\,$\to$\,\texttt{<mem>}/\texttt{<mem-idx>}, neutralized calls\,$\to$
\texttt{<fn>}), the same vocabulary-reduction principle used by
FastSpec~\cite{fastspec}. Two encoders are trained on the training corpus only:
a GPU-free static embedding (positive pointwise mutual information over a
windowed co-occurrence matrix followed by truncated SVD, in the spirit of
implicit-matrix-factorization word embeddings~\cite{levy2014}), and a contextual
masked-language-model Transformer~\cite{devlin2019} that treats each normalized
instruction as one token and yields per-instruction embeddings (an
assembly analogue of instruction embeddings~\cite{ding2019asm2vec}).

Table~\ref{tab:feat-ablation} isolates the representation with a RandomForest on
the locked test set. Purely learned sequence features reach 91\% accuracy with
zero hand-design; adding a small spec-derived structural descriptor recovers
most of the remaining gap; and hand features concatenated with the contextual
embedding exceed the hand baseline on macro-F1 (the gain concentrates on
minority classes). The static and contextual encoders behave differently: the
contextual embedding is \emph{complementary} to the hand features, whereas the
static one is largely redundant.

\begin{table}[t]
\centering
\caption{Feature representation ablation (RandomForest, locked test set,
  sequence-level mean-pooled features). The PPMI+SVD and MLM rows use only
  self-supervised embeddings; ``$+$\,spec-structural'' adds a spec-derived
  opcode-category histogram (not hand-tuned); the last row concatenates the
  58 hand features with the contextual embedding.}
\label{tab:feat-ablation}
\footnotesize
\begin{tabularx}{\columnwidth}{@{}Xrr@{}}
\toprule
Feature set & Test acc. & Macro-F1 \\
\midrule
Hand-engineered (58)              & 95.2\% & 80.3 \\
Learned, static PPMI+SVD (64)     & 91.3\% & 75.4 \\
Learned + spec-structural (84)    & 93.7\% & 78.6 \\
Hand $+$ contextual MLM (122)     & \textbf{95.5\%} & \textbf{89.3} \\
\bottomrule
\end{tabularx}
\end{table}

The decisive test replaces \emph{only the per-node feature vector} inside the
full GINE model (Section~\ref{sec:model}), aligning one contextual embedding to
each PDG node. Every other component is held fixed and identical across all
arms: the same spec-derived typed edges, the same 58 inline features
(Section~\ref{sec:inline}), and the same global and architecture features.
Table~\ref{tab:gine-parity} reports the single-seed outcome on the locked test
set that first suggested parity (a 0.6\,pp gap). That reading does not survive
scrutiny. A single seed cannot establish equivalence, and \emph{absence of a
significant difference is not evidence of equivalence}. We therefore ran a
pre-registered two-one-sided-test (TOST) equivalence analysis over ten seeds
per arm (margin $0.5$\,pp accuracy) on a compact GINE proxy, with both the
original and a scale-tuned encoder. In every case the learned node vector is
\emph{significantly lower} than the hand vector ($-2.5$ to $-3.0$\,pp accuracy;
$90\%$ CI entirely below zero; neither equivalent nor non-inferior). To rule out
the proxy itself as the cause, we repeated the analysis with the \emph{full}
production GINE (five seeds per arm, per-ISA spec builder, scale-tuned encoder;
Table~\ref{tab:full-tost}) and obtained the same verdict at the same effect size
($-2.58$\,pp, $90\%$ CI $[-3.81,-1.35]$). Concatenating the two (``both'') does
not exceed hand either ($-2.42$\,pp)---at the node level, once
the typed-edge graph structure is present, an added learned node representation
is redundant. We therefore \emph{withdraw the ``replaceable'' claim}: the
self-supervised encoder does not replace the hand-designed node features inside
the graph model. What it \emph{does} provide is complementarity at the sequence
level: a mean-pooled MLM embedding concatenated with the inline features raises
macro-F1 by $+10.6$\,pp over the hand features alone (RandomForest,
Table~\ref{tab:feat-ablation}), because sequence pooling discards the graph
structure that the GINE edges already supply, so the contextual encoder
contributes orthogonal signal there. The honest headline is thus
\emph{complementary learned features on minority classes}, not automated
replacement of the node vector.

\begin{table}[t]
\centering
\caption{Node-feature ablation in the full GINE model (locked test set,
  \emph{single seed} per row, identical training). \emph{Only the per-node vector
  varies}; the spec-derived typed edges, the 58 inline features, and the
  global/architecture features are present and identical in every row. The small
  single-seed gap does not hold up: a ten-seed TOST (main text) shows the learned
  node vector is significantly lower and not equivalent to hand. Reported here
  only to document the integration; not evidence of parity.}
\label{tab:gine-parity}
\footnotesize
\begin{tabularx}{\columnwidth}{@{}Xrrr@{}}
\toprule
Per-node vector & dim & Test acc. & Macro-F1 \\
\midrule
Hand nodes (reference)          & 41  & 95.6\% & 81.0 \\
Learned nodes                   & 65  & 95.0\% & 80.1 \\
Hand $+$ learned nodes          & 105 & 94.7\% & 78.0 \\
\bottomrule
\end{tabularx}
\end{table}

\begin{table}[t]
\centering
\caption{Multi-seed TOST in the \emph{full} production GINE (five seeds per arm:
  $\{42,1,7,13,21\}$; per-ISA spec builder; scale-tuned encoder; identical
  training budget). Differences are against the hand arm at an equivalence margin
  of $0.5$\,pp accuracy. Both learned arms are significantly \emph{worse}, with
  the $90\%$ CI entirely below zero---the compact-proxy verdict was not a proxy
  artifact. Note also the width of the hand arm's interval: the best single run
  reached $97.07\%$, which is top-of-range rather than typical.}
\label{tab:full-tost}
\footnotesize
\begin{tabularx}{\columnwidth}{@{}Xrrl@{}}
\toprule
Per-node vector & Test acc. (mean $\pm$ 90\% CI) & Macro-F1 & TOST verdict \\
\midrule
Hand (reference) & $\mathbf{96.14 \pm 1.59}$ & 85.6 & --- \\
Learned          & $93.57 \pm 0.60$ & \textbf{86.2} & not equivalent ($-2.58$) \\
Hand $+$ learned & $93.72 \pm 0.82$ & 84.8 & not equivalent ($-2.42$) \\
\bottomrule
\end{tabularx}
\end{table}

Several bounds keep this result honest. First, only the \emph{node-feature}
component is externalized; the nine edge types and the 58 inline features
(Section~\ref{sec:inline}) remain hand-specified and are used unchanged, so the
load-bearing attack knowledge (the speculative edge topology and windows) is
hand-authored in every arm. An edge-ablation over five seeds quantifies its
role: removing the speculative edge types (\textsc{spec\_conditional/indirect/
return}, \textsc{rsb\_chain}, \textsc{cache\_temporal}) costs $\approx 0$\,pp
accuracy but $-3.75$\,pp macro-F1, so the taxonomy matters for minority classes,
not headline accuracy (which rides node-level opcode statistics). Second, the
addressing grammar and all edge windows are now sourced from the specification
(the tokenizer carries no ISA-literal regex; the RSB window is spec-driven rather
than a module constant), but the speculation-primitive \emph{tags}
(\texttt{is\_secret\_source}, \texttt{is\_transmitter}, \dots) remain
hand-authored security semantics, and the \texttt{base} document still carries
the union of the x86 and ARM grammars; a genuinely new ISA such as RISC-V is
expressible in the schema but not yet exercised, so ``a new ISA needs only a spec
file'' is untested extrapolation. Third, and most importantly, the accuracy
numbers above use a \emph{random} record-level split, which leaks augmented
variants of the same base gadget across train and test. Under a group-hold-out
(no base gadget on both sides) accuracy falls $4.9$\,pp and macro-F1 $7.9$\,pp;
under an optimization-level hold-out (train O0--O2/Os, test O3) macro-F1
collapses from $94$ to $39$ while accuracy stays high only because \textsc{benign}
dominates. Cross-optimization generalization for the attack classes is therefore
weak, and the locked random test overstates it. Within these bounds, graph
construction is fully spec-driven, but ``automated feature engineering at no
accuracy cost'' overstates the case: the honest claims are a validated,
drift-proof, spec-driven front end and complementary learned features on
minority classes.

\subsection{Class- and Architecture-Conditioned Generation}
\label{sec:generation}

The same normalized token vocabulary supports generation. A decoder-only
Transformer~\cite{vaswani2017} is trained with next-token prediction on
sequences prefixed by two control tokens, a target class and a target
architecture:
\[
  [\,\langle\textsc{cls}\rangle_c,\ \langle\textsc{arch}\rangle_a,\
     i_1, i_2, \ldots, \langle\textsc{eos}\rangle\,],
\]
so decoding from $[\langle\textsc{cls}\rangle_c,\langle\textsc{arch}\rangle_a]$
conditions the sample on both. Because every emitted token is a normalized
instruction, output is grammar-valid by construction; a spec-driven
\emph{realizer} then instantiates operand placeholders with concrete registers,
immediates (biased toward attack-critical constants such as the page and
cache-line sizes), and addressing forms drawn from the target ISA's
specification, yielding assembly that the PDG builder parses.

Conditioning is verified independently: for each (class, architecture) we sample
gadgets and ask a reference classifier---the contextual encoder of
Section~\ref{sec:learned-features} with a RandomForest fit on \emph{real}
data---which class each sample resembles. The class hit-rate is compared to the
class prior, and an \emph{ISA-purity} score reports the fraction of
ISA-decisive opcodes native to the requested architecture.
Table~\ref{tab:gen-results} summarizes the outcome. Generation is strongly
class-conditioned ($5.7\times$ mean lift over the prior) and almost perfectly
architecture-consistent ($\geq 96\%$ purity); generated gadgets are largely
novel (not present in training) and exhibit the expected structure---for
Spectre-V1, a bounds compare, a conditional branch, and a secret-indexed load.
Per-class hit-rate tracks the training corpus's (class, architecture) coverage:
strong across x86, weaker for attack classes that were compiled predominantly
for x86 and are thus thin on ARM, which motivates broadening the compile matrix
rather than any change to the method.

\begin{table}[t]
\centering
\caption{Conditioned generation, verified by an independent reference
  classifier. Mean over classes; lift is hit-rate over class prior.}
\label{tab:gen-results}
\footnotesize
\begin{tabularx}{\columnwidth}{@{}Xrr@{}}
\toprule
Conditioning & Class lift & ISA-purity \\
\midrule
Class only                       & $7.4\times$ & --- \\
Class $+$ architecture (x86\_64) & \multirow{2}{*}{$5.7\times$} & 97.6\% \\
Class $+$ architecture (arm64)   &                              & 96.1\% \\
\bottomrule
\end{tabularx}
\end{table}

This conditioned, grammar-valid, architecture-aware generator is the candidate
source for the discovery loop: its samples are ranked by a learned surrogate and
the top candidates sent to microarchitectural simulation for confirmation and
minimization---the steps that turn recognition into discovery.

\section{Conclusions}
\label{sec:conclusion}

We presented the full methodology of \textsc{SpecExec}, a system that achieves
\textbf{95.75\%} accuracy (v54) on a 10-class speculative execution
vulnerability detection task directly from compiled assembly, with a weighted
F1 of 0.958 and macro recall of 0.920 (Table~\ref{tab:results-full}).
Key contributions are:

\begin{enumerate}[itemsep=1pt]
  \item A 9-edge-type PDG (Section~\ref{sec:pdg}) that explicitly encodes
    speculative execution paths; SPEC\_RETURN consistently receives the
    highest learned scale ($s\approx1.25$), confirming return-based edges
    as primary discriminating structure (Table~\ref{tab:edge-scales}).
  \item 58-dimensional inline features (Section~\ref{sec:inline}), including
    a Meltdown page-probe feature (\texttt{has\_page\_probe\_load}) that
    resolves L1TF/MDS confusion using the $2^{12}$-byte page stride as a
    discriminating signature. L1TF recall reaches 81.1\% (v54) vs.\ a
    baseline where it was confused with MDS.
  \item A provably correct augmentation suite (9 transforms, 143 tests,
    Section~\ref{sec:augmentation}) that found and fixed 6 implementation
    bugs; the corrected augmentation improves SPECTRE\_V2 recall by
    +11.7\,pp and RETBLEED by +6.6\,pp (v55 vs.\ v54).
  \item A locked test set enabling version-controlled regression tracking;
    Figures~\ref{fig:cm} and~\ref{fig:training} show confusion matrices
    and training curves for direct visual comparison.
  \item An automation path (Section~\ref{sec:specdiscover}) that externalizes the
    hand-authored feature engineering behind a declarative ISA specification: it
    drives graph construction with zero divergence from the reference builder
    (Table~\ref{tab:spec-equiv}) and $98.80\%$ agreement with an independent
    assembler/disassembler oracle (274 triaged defects); a self-supervised
    encoder that, by a multi-seed equivalence test, does not replace but is
    \emph{complementary} to the hand-designed node features ($+10.6$\,pp macro-F1
    at the sequence level); and a class- and architecture-conditioned generator
    that produces novel, grammar-valid gadgets (Table~\ref{tab:gen-results}),
    establishing the generative front end for simulator-in-the-loop gadget
    discovery.
  \item An honest evaluation of generalization: leakage-controlled (group- and
    optimization-level) hold-outs show the random-split accuracy is optimistic
    and cross-optimization minority-class macro-F1 is weak---so the classifier is
    best framed as \emph{gadget-candidate triage} feeding principled
    microarchitectural verification, not standalone vulnerability adjudication.
\end{enumerate}

\begin{thebibliography}{9}

\bibitem{safeside}
Google Project Zero, \textit{SafeSide: Hardware-confirmed speculative
execution proof-of-concepts.}
\url{https://github.com/google/safeside}, 2019.

\bibitem{spectector}
M.\ Guarnieri, B.\ K\"opf, J.\ F.\ Morales, J.\ Reineke, and A.\ S\'anchez,
\textit{Spectector: Principled Detection of Speculative Information Flows.}
IEEE~S\&P, 2020.

\bibitem{kocher1998}
P.\ Kocher et al.,
\textit{Spectre Attacks: Exploiting Speculative Execution.}
IEEE~S\&P, 2019.

\bibitem{fastspec}
Y.\ Wang et al.,
\textit{FastSpec: Scalable Generation and Detection of Spectre Gadgets using
Neural Embeddings.}
EuroS\&P, 2021.

\bibitem{hu2020strategies}
W.\ Hu, B.\ Liu, J.\ Gomes, M.\ Zitnik, P.\ Liang, V.\ Pande, and J.\ Leskovec,
\textit{Strategies for Pre-training Graph Neural Networks.}
ICLR, 2020.

\bibitem{xu2018jk}
K.\ Xu, C.\ Li, Y.\ Tian, T.\ Sonobe, K.\ Kawarabayashi, and S.\ Jegelka,
\textit{Representation Learning on Graphs with Jumping Knowledge Networks.}
ICML, 2018.

\bibitem{khosla2020supervised}
P.\ Khosla et al.,
\textit{Supervised Contrastive Learning.}
NeurIPS, 2020.

\bibitem{loshchilov2019decoupled}
I.\ Loshchilov and F.\ Hutter,
\textit{Decoupled Weight Decay Regularization.}
ICLR, 2019.

\bibitem{wong2011}
B.\ Wong,
\textit{Points of view: Color blindness.}
Nature Methods, 8(6):441, 2011.

\bibitem{levy2014}
O.\ Levy and Y.\ Goldberg,
\textit{Neural Word Embedding as Implicit Matrix Factorization.}
NeurIPS, 2014.

\bibitem{devlin2019}
J.\ Devlin, M.-W.\ Chang, K.\ Lee, and K.\ Toutanova,
\textit{BERT: Pre-training of Deep Bidirectional Transformers for Language
Understanding.}
NAACL, 2019.

\bibitem{ding2019asm2vec}
S.\ H.\ H.\ Ding, B.\ C.\ M.\ Fung, and P.\ Charland,
\textit{Asm2Vec: Boosting Static Representation Robustness for Binary Clone
Search against Code Obfuscation and Compiler Optimization.}
IEEE~S\&P, 2019.

\bibitem{vaswani2017}
A.\ Vaswani et al.,
\textit{Attention Is All You Need.}
NeurIPS, 2017.

\end{thebibliography}

\end{document}
